\chapter{Deep Dive: Observability --- Traces Across Feign and Kafka, Sampling, and Alerts That Matter}
\label{ch:dd-observability}

Chapter~24 taught you to debug the stack from the outside in ---
\code{describe} before \code{logs} --- and Chapter~27 named the trio
that comes next: metrics, traces, alerts. This deep dive opens the
wiring that is already in every service and asks what it actually does.
The answer is less than the configuration suggests: the tracing
dependencies are on every classpath, the log pattern prints trace ids,
the cluster samples at zero, and the Prometheus endpoint that Chapter~27
counts on does not exist yet. Reading why is the fastest route to
understanding how distributed tracing works at all.

Three questions organize the chapter:

\begin{enumerate}
  \item \textbf{Where does a trace start, and how does it survive} the
  hops browser $\rightarrow$ gateway $\rightarrow$ Feign $\rightarrow$
  Kafka $\rightarrow$ consumer --- and where in this project does it
  quietly die?
  \item \textbf{What does sampling cost} at 1.0, at 0.0, and at the
  numbers production systems really use?
  \item \textbf{Which signals deserve an alert}, and what has to be
  added before any alert can fire?
\end{enumerate}

\section{Why this, and what it solves}

Logs answer ``what did this process say?'' Metrics answer ``how is the
system doing, in aggregate, over time?'' Traces answer the question the
other two cannot: ``what happened to \emph{this one request} as it
crossed five services and a message broker?'' In a monolith a stack
trace is a trace. In this project a single ``create assignment'' click
touches the gateway, course-service, auth-service (over Feign),
Kafka, and notification-service; no single log file sees all of it.

The project chose Micrometer Tracing with the Brave bridge, reporting to
Zipkin. Each service's \code{pom.xml} carries the same pair:

\begin{lstlisting}[language=yaml]
<dependency>
  <groupId>io.micrometer</groupId>
  <artifactId>micrometer-tracing-bridge-brave</artifactId>   (*@\co{1}@*)
</dependency>
<dependency>
  <groupId>io.zipkin.reporter2</groupId>
  <artifactId>zipkin-reporter-brave</artifactId>             (*@\co{2}@*)
</dependency>
\end{lstlisting}

\begin{description}
  \item[\co{1}] Micrometer Tracing is the \emph{facade} Spring Boot~3
  instruments against; Brave (the Zipkin project's tracer) is one of
  two implementations, OpenTelemetry being the other. The bridge is
  what turns Spring's \emph{observations} --- the Micrometer Observation
  API that every Boot~3 starter emits --- into spans.
  \item[\co{2}] The reporter ships finished spans to Zipkin over HTTP.
  Swapping this one artifact for \code{opentelemetry-exporter-otlp}
  points the same spans at Tempo, Jaeger, or a vendor.
\end{description}

\begin{decision}[Brave and Zipkin, not OpenTelemetry]
Brave plus Zipkin is the shortest path from zero to a working trace
view: one container, no collector, no protocol negotiation, and Spring
Boot's auto-configuration knows both intimately. That is the right
choice for learning the mechanics and for a homelab. What would flip
it: a second language in the stack (OpenTelemetry's SDKs cover all of
them with one wire protocol), a vendor backend (Datadog, Honeycomb,
Grafana Cloud all speak OTLP natively), or the need for tail-based
sampling, which requires a collector between the apps and the store.
The migration is mechanical --- swap two dependencies, change the
endpoint property --- because the code never touched Brave directly.
Interviewers ask ``Zipkin or Jaeger or Tempo?'' expecting exactly this
answer: the backend is replaceable, the instrumentation standard is the
decision that lasts.
\end{decision}

\section{How a trace travels through this project}

\subsection{The configuration, in every service}

The five service documents in config-server all end the same way:

\begin{lstlisting}[language=yaml]
management:
  tracing:
    sampling:
      probability: 1.0                                   (*@\co{1}@*)
  zipkin:
    tracing:
      endpoint: http://localhost:9411/api/v2/spans       (*@\co{2}@*)
  endpoints:
    web:
      exposure:
        include: health,info,metrics,prometheus          (*@\co{3}@*)
logging:
  pattern:
    level: "%5p [${spring.application.name:},%X{traceId:-},%X{spanId:-}]"  (*@\co{4}@*)
\end{lstlisting}

\begin{description}
  \item[\co{1}] Every request is traced. Fine on a laptop; see
  Section~\ref{sec:sampling} for what it costs elsewhere. In the
  cluster, every Deployment overrides this to \code{0.0} with the
  comment ``Zipkin is not in the cluster yet''.
  \item[\co{2}] Compose overrides the host to \code{zipkin:9411}; the
  reporter POSTs JSON batches here.
  \item[\co{3}] Exposure is a \emph{list of names}. Naming
  \code{prometheus} does not create the endpoint; the endpoint exists
  only when \code{micrometer-registry-prometheus} is on the classpath,
  and no \code{pom.xml} in the repository has it. Today
  \code{/actuator/prometheus} answers 404 on every service --- the
  Chapter~27 plan assumes an endpoint that still has to be added.
  \item[\co{4}] The log pattern reads two MDC keys that the tracing
  bridge populates. With no active span the fallbacks \code{-} print
  empty, which is itself diagnostic: an empty trace id in a log line
  means the code ran outside any traced context.
\end{description}

\subsection{Hop by hop}

Follow one \code{POST /api/courses/7/assignments} from the browser.

\textbf{Browser $\rightarrow$ gateway.} The browser sends no trace
header; \code{axios} knows nothing of tracing. The gateway is a WebFlux
application and Spring Cloud Gateway is instrumented through the
Observation API, so the server-side observation of the incoming request
\emph{starts} a new trace: a random 128-bit trace id and a root span.
When the gateway forwards to course-service, the reactive
\code{WebClient} it uses is instrumented too, and the propagator writes
the context into the outgoing request. Spring Boot~3 defaults matter
here: Brave is configured to \emph{produce} W3C Trace Context
(\code{traceparent: 00-<traceId>-<spanId>-01}) and to \emph{consume}
W3C \emph{and} both B3 formats. So this project speaks W3C on the wire
and would still understand an older B3 caller.

\textbf{Gateway $\rightarrow$ course-service.} Course-service is servlet
MVC. Its server observation reads \code{traceparent}, continues the
same trace, and opens a child span. Everything logged inside the
controller now carries that trace id --- the pattern in \co{4} above
prints it.

\textbf{course-service $\rightarrow$ auth-service over Feign.}
\code{AuthServiceClient} is a Spring Cloud OpenFeign client. OpenFeign
ships a Micrometer capability, enabled by default, that wraps each call
in an observation; the Brave propagator injects \code{traceparent} into
the Feign request. Auth-service continues the trace. In Zipkin the Feign
call appears as a client span in course-service nested under the
controller span, with a matching server span in auth-service. If
Resilience4j's circuit breaker short-circuits the call
(Chapter~6), no client span is created --- the fallback ran, the wire
was never touched, and the trace shows the gap.

\textbf{course-service $\rightarrow$ Kafka.} Here the trace dies. Read
the project's producer:

\begin{lstlisting}[language=Java]
@Bean
public ProducerFactory<String, Object> producerFactory() {
    Map<String, Object> configProps = new HashMap<>();
    configProps.put(BOOTSTRAP_SERVERS_CONFIG, bootstrapServers);
    configProps.put(KEY_SERIALIZER_CLASS_CONFIG, StringSerializer.class);
    configProps.put(VALUE_SERIALIZER_CLASS_CONFIG, JsonSerializer.class);
    return new DefaultKafkaProducerFactory<>(configProps);
}

@Bean
public KafkaTemplate<String, Object> kafkaTemplate() {
    return new KafkaTemplate<>(producerFactory());       (*@\co{1}@*)
}
\end{lstlisting}

\begin{description}
  \item[\co{1}] A hand-built \code{KafkaTemplate}. Spring Kafka
  \emph{can} propagate trace context --- it writes \code{traceparent}
  into the record's headers and starts a producer span --- but only
  when \code{setObservationEnabled(true)} is called on the template, or
  \code{spring.kafka.template.observation-enabled=true} is set for the
  auto-configured one. This bean does neither, and the property is
  absent from every config document. The record leaves without a
  header. On the consuming side, \code{@KafkaListener} in
  notification-service would need
  \code{spring.kafka.listener.observation-enabled=true} to read one.
  Neither half is on. The e-mail that results from the click is, to
  Zipkin, an unrelated event.
\end{description}

Auth-service's \code{UserRegisteredPublisher} uses the auto-configured
template and has the same gap for the same reason: the property, not
the bean, is what is missing.

\textbf{The scheduler boundary.} \code{CourseLifecycleScheduler} runs
\code{@Scheduled} at 01:00. Nothing instruments a cron trigger; its log
line ``Auto-archived 3 course(s)'' prints with empty trace and span ids.
The same is true of any \code{@Async} method unless the executor is
wrapped with a context-propagating decorator. A trace is a property of a
\emph{request}; work that no request started has no trace unless you
give it one.

\subsection{Correlating logs}

With the pattern from \co{4}, a grep across all five services' logs for
one trace id assembles the story Zipkin draws --- useful precisely when
Zipkin is down or sampled out:

\begin{lstlisting}[language=bash]
$ docker compose logs --no-color api-gateway course-service \
    auth-service | grep 6f2c1a9e4b3d8e10
api-gateway     | INFO [api-gateway,6f2c1a9e4b3d8e10,9a1b...] ...
course-service  | INFO [course-service,6f2c1a9e4b3d8e10,c04d...] ...
auth-service    | INFO [auth-service,6f2c1a9e4b3d8e10,e77f...] ...
\end{lstlisting}

Notice what is absent: no \code{notification-service} line, because no
trace id crossed Kafka. The grep shows the hole exactly as Zipkin would.

\section{Sampling}
\label{sec:sampling}

\code{probability: 1.0} means every trace is recorded and every span
shipped. The cost is linear in traffic: each span is a few hundred bytes
of JSON, batched and POSTed by a background thread, then indexed and
stored by Zipkin. At the project's traffic that is nothing. At
1{,}000 requests per second with six spans each it is 6{,}000 spans per
second, roughly 2\,MB/s of ingest and tens of gigabytes of storage a
day --- most of it recording requests that succeeded in 20\,ms and
that nobody will ever look at.

\emph{Head-based} sampling --- what \code{probability} configures ---
decides at the root span whether to keep the trace, and propagates the
decision in the \code{traceparent} flags so every downstream service
agrees. It is cheap and consistent, and it is blind: at 0.01 it keeps
1\% of traces chosen at random, so the one request that failed at
3\,a.m. was kept with probability 0.01. \emph{Tail-based} sampling
defers the decision until the whole trace is known and keeps every
trace that errored or exceeded a latency budget, plus a random sample
of the rest. It needs a component that buffers whole traces before the
store --- an OpenTelemetry Collector --- which is why it is out of
reach of the Brave-to-Zipkin direct wiring.

\begin{decision}[the sampling rate]
Laptop and Compose: \code{1.0}, because you are debugging and traffic
is you. Cluster today: \code{0.0}, honest about the absence of a
Zipkin pod --- sampling at zero also disables span creation cost, not
only shipping. Cluster with Zipkin: start at \code{0.1}; drop to
\code{0.01} when the store grows faster than you read it. The senior
answer to ``what sample rate?'' is ``head-based 1--10\% for the
baseline, tail-based for errors and slow requests once there is a
collector'' --- and the follow-up is always ``and error traces are
what you actually want, so budget for the collector''.
\end{decision}

\section{Metrics: what exists and what does not}

Actuator's \code{/actuator/metrics} already exists on every service and
lists a few hundred meters: JVM memory and GC, HTTP server requests
tagged by \code{uri}, \code{status} and \code{method}, Kafka client
metrics, HikariCP pool state, and Resilience4j circuit-breaker state
in course-service. It answers one meter at a time, by name, over HTTP
--- a debugging tool, not a monitoring one. Monitoring means a system
that \emph{pulls} all of them every 15 seconds, stores them as time
series, and evaluates rules. That system is Prometheus, and its
endpoint is missing (see \co{3} above).

Two vocabularies to carry into interviews. \textbf{RED} is for
request-driven services: Rate, Errors, Duration --- the three the
\code{http.server.requests} timer already provides. \textbf{USE} is
for resources: Utilization, Saturation, Errors --- CPU, heap, the
Hikari pool, Kafka consumer lag. A dashboard per service should show
RED at the top and USE below it; an alert should fire on a RED symptom
the user feels, and USE panels explain why.

\section{Pros and cons}

\begin{center}
\begin{tabular}{@{}p{0.47\linewidth}p{0.47\linewidth}@{}}
\toprule
\textbf{Strengths} & \textbf{Weak points, today} \\
\midrule
Zero code changes: every hop is instrumented by starters, and the log
pattern is one line. &
Kafka propagation is off in both producer and consumer; the most
interesting path in the system is the one that is invisible. \\
Vendor-neutral facade: Brave and Zipkin are swappable for OTLP
without touching source. &
No Prometheus registry on any classpath; \code{prometheus} in the
exposure list is aspirational. \\
W3C \code{traceparent} on the wire, B3 accepted --- interoperable
with anything modern. &
Head-based sampling only; an error at 1\% sampling is kept 1\% of the
time. \\
Log correlation works even without Zipkin. &
Zipkin runs \code{:latest} with no volume: traces vanish on every
\code{compose down}. \\
\bottomrule
\end{tabular}
\end{center}

\section{What breaks}

\begin{pitfall}[the trace that ends at Kafka]
Symptom: in Zipkin, the \code{POST /assignments} trace ends with the
controller span; the notification that followed is a separate one-span
trace, or absent. Cause: neither \code{spring.kafka.template.\allowbreak
observation-enabled} nor \code{spring.kafka.listener.\allowbreak
observation-enabled} is set, so no \code{traceparent} header is written
to or read from the record. Fix: both properties, in both services'
config documents (course-service and auth-service produce;
notification-service consumes) --- and the hand-built template in
\code{KafkaProducerConfig} must call \code{setObservationEnabled(true)}
because properties only configure the auto-configured bean. The
Section~\ref{sec:enhance} listing does all three.
\end{pitfall}

\begin{pitfall}[high-cardinality tags]
Symptom: a service's heap climbs for days and \code{/actuator/metrics}
takes seconds to answer; eventually \code{OOMKilled}. Cause: a tag whose
value set is unbounded. The classic is a handler mapped as
\code{/courses/\{id\}} that someone rewrote to build the path by hand
--- the \code{uri} tag then carries \code{/courses/42},
\code{/courses/43}, \ldots, one time series each, forever. The same
happens when a span is tagged with a user id or an e-mail. Fix: tag
with the \emph{template} (\code{/courses/\{id\}}), never the value;
put identifiers in the span's log annotations or the log line, which
are not indexed by value. Spring's default \code{uri} tag is the
template precisely to avoid this.
\end{pitfall}

\begin{pitfall}[a 401 that has no trace]
Symptom: a user reports ``I got logged out'' and there is no trace for
their request. Cause: the gateway's \code{JwtAuthenticationFilter}
rejects at order \code{-1}, before routing; the server observation
still records the request, but with the sampling at 0.0 in the cluster
nothing is exported anywhere. Fix is not tracing at all: this is a
\emph{metric} question. Count \code{http.server.requests} with
\code{status=401} at the gateway and alert on a jump --- one bad
deploy of \code{JWT\_SECRET} shows as a cliff on that counter within
15 seconds.
\end{pitfall}

\section{What breaks at 3\,a.m.}

Zipkin dies at 3\,a.m.; the services keep running.

\begin{itemize}
  \item \textbf{Request latency: unchanged.} The reporter is
  asynchronous. Spans finish into an in-memory queue; a background
  thread drains it in batches over HTTP. The request thread never waits
  on Zipkin. This is the property to state first in an interview,
  because the fear ``tracing slows my service'' is the common objection.
  \item \textbf{Memory: bounded.} The queue has a fixed capacity
  (\code{queuedMaxSpans}, 10{,}000 by default). When it is full, new
  spans are \emph{dropped}, counted, and a warning is logged --- the
  reporter does not grow the heap to keep them. You lose visibility,
  not the service.
  \item \textbf{Logs: still correlated.} Trace ids are generated
  locally; the MDC pattern keeps printing them. The grep in the
  correlation section works exactly as before. This is the argument for
  always having the log pattern, regardless of the trace backend.
  \item \textbf{Recovery: automatic.} The reporter retries each batch on
  the next drain; when Zipkin returns, new spans flow. The spans dropped
  during the outage are gone --- there is no replay.
\end{itemize}

Now the other direction: \emph{traffic} spikes at 3\,a.m. with sampling
at 1.0. The queue fills, spans drop, and the warning log itself becomes
noise at hundreds of lines per second. That is the concrete form of the
sampling trade-off --- and the reason to set the rate before the spike,
not during it.

\section{Enhancing the resolution}
\label{sec:enhance}

Four additions turn the wiring into a monitoring system. Each is
complete; add them in order.

\textbf{1. The Prometheus endpoint.} One dependency per service (the
parent's Boot BOM manages the version):

\begin{lstlisting}[language=yaml]
<dependency>
  <groupId>io.micrometer</groupId>
  <artifactId>micrometer-registry-prometheus</artifactId>
</dependency>
\end{lstlisting}

With it on the classpath, the existing exposure list makes
\code{/actuator/prometheus} real. Add the tracing switches that are
missing, and a common tag so every series says which service it came
from:

\begin{lstlisting}[language=yaml]
management:
  prometheus:
    metrics:
      export:
        enabled: true
  metrics:
    tags:
      application: ${spring.application.name}   # every series, one tag
  tracing:
    sampling:
      probability: ${TRACING_SAMPLE:0.1}          # env-overridable
spring:
  kafka:
    template:
      observation-enabled: true    # producer span + traceparent header
    listener:
      observation-enabled: true    # consumer span, continues the trace
\end{lstlisting}

And the hand-built template in course-service, which ignores the
property:

\begin{lstlisting}[language=Java]
@Bean
public KafkaTemplate<String, Object> kafkaTemplate() {
    KafkaTemplate<String, Object> template =
            new KafkaTemplate<>(producerFactory());
    // Properties configure only the auto-configured template; a
    // hand-built one must opt in explicitly or traces end here.
    template.setObservationEnabled(true);
    return template;
}
\end{lstlisting}

\textbf{2. Scraping in the cluster.} With \code{kube-prometheus-stack}
installed (Chapter~27), a \code{ServiceMonitor} tells its operator
which Services to scrape. One object covers all five, by label:

\begin{lstlisting}[language=yaml]
apiVersion: monitoring.coreos.com/v1
kind: ServiceMonitor
metadata:
  name: ts-services
  namespace: ts
  labels:
    release: kube-prometheus-stack   # the selector the chart installs
spec:
  selector:
    matchLabels:
      monitoring: actuator           # add this label to each Service
  endpoints:
    - port: http                     # the Service port *name*
      path: /actuator/prometheus
      interval: 15s
\end{lstlisting}

The auth-service and course-service Services must permit
\code{/actuator/prometheus} in their SecurityConfig for this to work
--- the same permit that Chapter~17 needed for probes. Restrict it to
the pod network with a NetworkPolicy rather than leaving it open to the
Ingress.

\textbf{3. Structured logs.} A grep across containers works on a
laptop; in the cluster, Loki or Elasticsearch parse JSON far more
reliably than a bracketed pattern. Logback's JSON encoder is a
dependency plus a \code{logback-spring.xml}:

\begin{lstlisting}[language=yaml]
<dependency>
  <groupId>net.logstash.logback</groupId>
  <artifactId>logstash-logback-encoder</artifactId>
  <version>8.0</version>
</dependency>
\end{lstlisting}

\begin{lstlisting}
<configuration>
  <springProperty name="app" source="spring.application.name"/>
  <appender name="JSON"
            class="ch.qos.logback.core.ConsoleAppender">
    <encoder class="net.logstash.logback.encoder.LogstashEncoder">
      <customFields>{"application":"${app}"}</customFields>
      <!-- traceId and spanId come from the MDC automatically -->
      <includeMdcKeyName>traceId</includeMdcKeyName>
      <includeMdcKeyName>spanId</includeMdcKeyName>
    </encoder>
  </appender>
  <springProfile name="!local">
    <root level="INFO"><appender-ref ref="JSON"/></root>
  </springProfile>
  <springProfile name="local">
    <include resource="org/springframework/boot/logging/logback/base.xml"/>
  </springProfile>
</configuration>
\end{lstlisting}

The profile split keeps the human-readable pattern on a laptop and emits
one JSON object per line everywhere else. Every line now carries
\code{traceId} as a field; a log query for a trace id is an index
lookup, not a scan.

\textbf{4. Two alerts tied to what users feel.} Prometheus rules, in
the operator's \code{PrometheusRule} kind:

\begin{lstlisting}[language=yaml]
apiVersion: monitoring.coreos.com/v1
kind: PrometheusRule
metadata:
  name: ts-slo
  namespace: ts
  labels: { release: kube-prometheus-stack }
spec:
  groups:
    - name: ts.slo
      rules:
        - alert: GatewayErrorBudgetBurn
          expr: |
            sum(rate(http_server_requests_seconds_count{
                  application="api-gateway", status=~"5.."}[5m]))
            /
            sum(rate(http_server_requests_seconds_count{
                  application="api-gateway"}[5m])) > 0.02
          for: 10m
          labels: { severity: page }
          annotations:
            summary: "Gateway 5xx ratio above 2% for 10 minutes"
        - alert: NotificationConsumerLag
          expr: |
            sum by (topic) (
              kafka_consumer_fetch_manager_records_lag_max{
                application="notification-service"}) > 500
          for: 15m
          labels: { severity: ticket }
          annotations:
            summary: "notification-group is {{ $value }} records behind"
\end{lstlisting}

The first is a \emph{symptom} alert on the RED ``errors'' ratio at the
one place every request passes; 2\% over ten minutes is a starting
budget, not a law. The second is the Chapter~27 answer to ``which Kafka
metric first'': lag growing for 15 minutes means mail is not going out,
whatever the pods' health says. Note the meter name: Spring Kafka
registers the client's own \code{records-lag-max} under Micrometer's
naming, which the Prometheus registry renders with underscores.

\begin{handson}[prove the trace crosses Kafka]
After adding the two \code{observation-enabled} properties and the
template line, restart course-service and notification-service under
Compose, create an assignment through the UI, then ask Zipkin for the
trace by the id printed in course-service's log:
\begin{lstlisting}[language=bash]
$ curl -s "http://localhost:9411/api/v2/trace/6f2c1a9e4b3d8e10" \
    | jq -r '.[] | "\(.localEndpoint.serviceName)  \(.name)"'
api-gateway            http post
course-service         http post /courses/{id}/assignments
course-service         http get                   # Feign to auth
auth-service           http get /users/{id}
course-service         ts.assignment.created send
notification-service   ts.assignment.created process
\end{lstlisting}
Before the change the last two lines do not exist. The consumer span's
parent is the producer span; the e-mail is now attributable to the
click that caused it.
\end{handson}

\section{Outside the project}

Production stacks converge on the same shape with different names.
Instrumentation: OpenTelemetry SDKs or the Java agent, which
auto-instruments without even the starters. Collection: an
OpenTelemetry Collector per node or per cluster, which is where
tail-based sampling, redaction, and fan-out to several backends live.
Storage: Tempo for traces, Loki for logs, Prometheus or Mimir for
metrics --- the Grafana trio --- or Jaeger, Elasticsearch and Thanos;
or a vendor that takes OTLP and bills by volume. The interview thread
runs through all of them: ``how do you find a slow request across
services?'' is answered by trace ids in logs and a trace store; ``how
do you know before users tell you?'' by RED alerts on the edge; ``what
does tracing cost?'' by the sampling discussion above --- and the
senior version of every answer names the component that stays
(instrumentation and the trace id in every log line) and the ones that
are replaceable (everything behind the exporter).

\section{Questions from real interviews}

\subsection*{Q: p99 latency doubled last night. Walk me through finding the cause.}

Start with metrics, not logs: the \code{http\_server\_requests} histogram
at the gateway, split by \code{uri}, tells you in one query whether
\emph{every} route slowed (infrastructure --- node, network, GC) or one
route did (a code path or a downstream). Then pull one slow trace from
that window --- with head-based sampling at 1\% there will be a few
hundred --- and read the span waterfall: the slow span is either a
child (a downstream call, say the Feign hop to auth-service at 800\,ms
instead of 8) or a gap between children (the service itself, usually a
database call or a lock, which shows as time in the parent not covered
by any child). If it is a gap, the trace id takes you to that service's
logs for the exact minute; the JSON layout makes that an index lookup.
The senior detail is the order --- aggregate, then one exemplar, then
its logs --- and naming what you would have wished for: an exemplar
link from the histogram bucket straight to a trace, which
OpenTelemetry-native stacks provide and Zipkin does not.

\subsection*{Q: You cannot afford to trace 100\%. What do you sample, and what do you never drop?}

Head-based sampling at 1--10\% for the steady state, decided at the
gateway and propagated in the \code{traceparent} flags so downstream
services never disagree. Never drop errors and never drop the slow
tail --- which head-based sampling cannot promise, because the root
span is sampled before anyone knows the request will fail. So the real
answer has two stages: a collector doing tail-based sampling that keeps
every trace with a 5xx or above a latency budget, plus a 1\% random
baseline for the shape of normal traffic. Mention the cost honestly:
tail-based sampling buffers whole traces in the collector's memory for
the decision window, typically 10--30\,s, so the collector is sized by
traffic, not by what you keep. In this project that is the step from
\code{zipkin-reporter-brave} to an OTLP exporter plus a collector; the
application code does not change.

\subsection*{Q: A request disappears from the trace at Kafka. Why, and how do you fix it?}

Because trace context crosses Kafka only as record headers, and
nothing writes them unless the template is told to. In Spring Kafka
two properties do it, both false by default in Boot~3: on the
producer, \code{template.observation-enabled}; on the consumer,
\code{listener.observation-enabled} (both under \code{spring.kafka}).
The subtlety worth
saying out loud: properties only configure the auto-configured
template, so a hand-built \code{KafkaTemplate} bean --- exactly what
course-service's \code{KafkaProducerConfig} does --- needs
\code{setObservationEnabled(true)} in code. After the fix the consumer
span's parent is the producer span and the whole e-mail path hangs off
the original HTTP request. The follow-up an interviewer wants: the
consumer span is a \emph{child}, but it can run seconds or hours after
the request finished, so a trace view sorted by time will show a long
gap --- that gap is queue latency, and it is itself the metric to
watch.

\subsection*{Q: What would you page on for this stack?}

Symptoms users feel, expressed as SLOs, never raw resources. Two pages:
gateway 5xx ratio above the error budget for ten minutes --- the one
point every request passes --- and p99 latency at the gateway above
target for the same window. One ticket-level alert: notification
consumer lag growing for fifteen minutes, because that is ``mail is not
going out'' in the only form the system can express it. Everything
else --- CPU, heap, pod restarts, Hikari pool exhaustion --- is a
dashboard panel that explains a page, not a page itself; alerting on
CPU at 80\% wakes people for a system that is working. The senior
framing is the error budget: at a 99.9\% availability target you may
burn 43 minutes of errors a month, and the alert fires on the
\emph{burn rate}, which is why the ratio is measured over a window
rather than on a single failed request.

\subsection*{Q: One user's request touched five services. How do you find those five log lines?}

By the trace id, which every service prints because the log pattern
includes \code{\%X\{traceId\}} and the tracing bridge fills the MDC on
every hop. On a laptop that is a grep across container logs; in the
cluster it is a query on a log store that indexes the field, which is
why logs are shipped as JSON rather than parsed with a regex. Two
details separate a senior answer: first, the trace id must be
\emph{returned to the user} --- a \code{X-Trace-Id} response header
from the gateway --- so that a support ticket can quote it instead of
``around 3 p.m.''; second, anything that runs outside the request
(\code{@Scheduled}, \code{@Async}, a Kafka consumer without
propagation) logs an empty id, and you should know which of your code
paths those are before the incident.

\subsection*{Q: Zipkin is down. Does production slow down?}

No, and you should be able to say why from the reporter's design: spans
are handed to an in-memory queue and a background thread POSTs them in
batches, so no request thread ever waits on Zipkin; the queue is
bounded at 10{,}000 spans and drops on overflow, so memory is capped;
the failed batches are logged and retried on the next drain, so
recovery is automatic. What you lose is the spans from the outage
window --- there is no replay. The trap in the question is the
opposite case: a \emph{synchronous} exporter, or an unbounded queue,
would turn a monitoring outage into an application outage, which is
the property to check before adopting any agent or exporter. And the
mitigation for the lost spans is the same as for everything else in
this chapter: the trace id is still in every log line, so the logs
carry the correlation even when the trace store cannot.

\begin{quiz}
\begin{enumerate}
  \item Name the exact property or method call missing at each of the
  two Kafka-side halves of the project, and explain why setting the
  properties alone would fix auth-service but not course-service.
  \item A request arrives at the gateway carrying
  \code{X-B3-TraceId} from an old upstream proxy. Does the project
  continue that trace or start a new one, and which Boot~3 default
  decides it? What header leaves the gateway toward course-service?
  \item \code{/actuator/prometheus} is in every exposure list yet
  returns 404. Explain, and name the one artifact that changes the
  answer.
  \item Zipkin is down for an hour under normal traffic. State what
  happens to request latency, to heap, and to the spans produced
  during the hour --- and which of the three would change if the
  reporter were synchronous.
  \item At \code{probability: 0.01}, what fraction of \emph{failed}
  requests have a trace, and what architectural component is required
  to make that fraction 1.0 without tracing every success?
  \item A colleague adds \code{span.tag("email", user.getEmail())}
  to the login span. Predict the failure over the following weeks and
  give the rule that prevents it.
\end{enumerate}
\end{quiz}
